\documentclass[10pt]{beamer}
\usepackage[T1]{fontenc}   % \k{e} in Swiecki needs this
\usepackage{amsmath}
\usepackage{amssymb}
\usepackage{booktabs}
\usepackage{array}
\usepackage{graphicx}
\usepackage{xcolor}
\newcolumntype{L}[1]{>{\raggedright\arraybackslash}p{#1}}
%-------------------------------------------------------
%   Look
%-------------------------------------------------------
\usetheme{default}
\definecolor{accent}{RGB}{31,71,136}
\definecolor{dev}{RGB}{214,40,40}   % where we depart from the paper
\newcommand{\dev}[1]{\textcolor{dev}{#1}}
\setbeamercolor{structure}{fg=accent}
\setbeamercolor{frametitle}{fg=accent}
\setbeamerfont{frametitle}{series=\bfseries}
\setbeamerfont{framesubtitle}{size=\normalsize}
\setbeamertemplate{navigation symbols}{}
\setbeamertemplate{itemize items}[circle]
\setbeamertemplate{enumerate items}[default]
\setlength{\abovedisplayskip}{4pt}
\setlength{\belowdisplayskip}{4pt}
% Frame number, bottom right.
\setbeamertemplate{footline}{%
  \hfill{\usebeamercolor[fg]{structure}\scriptsize\insertframenumber}%
  \hspace{1.2em}\vspace{4pt}}
% One list spacing everywhere. Every itemize and enumerate opens with \itms.
\newcommand{\itms}{\itemsep5pt}
%-------------------------------------------------------
%   Frame helpers
%-------------------------------------------------------
% Where figures live. Each \figframe passes a bare file name and LaTeX searches
% these in order, so the same file builds on Overleaf (figures/ in the project)
% and in the repo (../analysis/results/figures/).
\graphicspath{{./}{figures/}{../analysis/results/figures/}}
% One figure per frame, identical sizing throughout.
\newcommand{\figframe}[3]{%
  \begin{frame}{#1}{#2}
    \vspace{6pt}
    \centering
    \includegraphics[width=\textwidth,height=0.76\textheight,keepaspectratio]{#3}
  \end{frame}}
% Appendix links. \ap{label} marks a word with a faint superscript that jumps to
% the appendix frame; \apback{label} is the small return arrow that goes back.
\newcommand{\ap}[1]{%
  \hyperlink{#1}{\structure{\textsuperscript{\tiny$\blacktriangleright$}}}}
\newcommand{\apback}[1]{%
  \vfill\raggedleft\hyperlink{#1}{\structure{\tiny$\blacktriangleleft$\,back}}}
% Placeholder frame: heading only, content to be added later
\newcommand{\todoframe}[3][TO ADD.]{%
  \begin{frame}{#2}{#3}
    \vfill
    \centering \textit{#1}
    \vfill
  \end{frame}}
% Divider frame at each section.
\AtBeginSection[]{%
  \begin{frame}[plain]
    \vfill
    \centering{\usebeamercolor[fg]{structure}\Large\insertsectionhead}
    \vfill
  \end{frame}}
\title{Title}
\author{Mitchell Ballif}
\date{\today}
\begin{document}
\begin{frame}[plain]
  \titlepage
\end{frame}
%=======================================================
\section{Introduction}
%=======================================================
\figframe{Introduction}{Motivation}{us_trade_motivation.pdf}
\begin{frame}{Introduction}{Literature}
  \structure{\textbf{Closed  economy}}
  \begin{itemize}\itms
    \item Technology story --- Baumol (1967); Ngai \& Pissarides (2007)
    \item Non-homothetic preferences --- Kongsamut, Rebelo \& Xie (2001);
          Boppart (2014); Comin, Lashkari \& Mestieri (2021)
  \end{itemize}
  \vspace{8pt}
  \structure{\textbf{Open economy}}
  \begin{itemize}\itms
    \item Comparative advantage --- Matsuyama (1992, 2009)
    \item Quantitative decompositions --- Uy, Yi \& Zhang (2013);
          \'Swi\k{e}cki (2017); Sposi (2019); Cravino \& Sotelo (2019);
          Sposi, Yi \& Zhang (2024)
  \end{itemize}
  \vspace{8pt}
  \structure{\textbf{Trade imbalances}}
  \begin{itemize}\itms
    \item Imbalances $\to$ sectoral composition --- Kehoe, Ruhl \& Steinberg (2018)
  \end{itemize}
\end{frame}
\begin{frame}{Introduction}{Research question}
  \vfill
  \begin{center}
  \large
  U.S.\ manufacturing's share of value added fell 13 percentage points
  between 1965 and 2014.\\[10pt]
  \structure{\textbf{How much of that decline was caused by falling trade costs?}}
  \end{center}
  \vfill
  \begin{center}
  \normalsize
  And how much of the answer depends on letting the trade deficit adjust?
  \end{center}
  \vfill
\end{frame}
\begin{frame}{Introduction}{Contribution}
  \vfill
  \begin{enumerate}\itms
    \item Extended baseline SC model (trade imbalance + financial frictions)
    \item Wrote math appendix (solving and inverting equilibrium conditions)
    \item Built long-run panel using World IO data
    \item Took the model to data (recovered shocks)
  \end{enumerate}
  \vfill
\end{frame}
%=======================================================
\section{Model}
%=======================================================
\begin{frame}{Model}{Environment}
  \begin{itemize}\itms
    \item Time is discrete, $t=1965,\dots,2014$.
    \item Two countries: US and the rest of the world,
          $i\in\mathcal{I}\equiv\{US,\,ROW\}$.
    \item Three sectors in each country, all tradable,
          $j\in\mathcal{J}\equiv\{a,m,s\}$.
    \item Representative household in each country:
      \begin{itemize}\itms
        \item Country $i$ endowed with $L_{i,t}$ units of homogeneous labor in
              every $t$ and $K_{i,0}$ units of homogeneous physical capital.
        \item Access to international financial markets by means of buying and
              selling one-period bonds with portfolio adjustment costs.
        \item Born with zero NFA, $B_{i,0}=0$.
      \end{itemize}
  \end{itemize}
\end{frame}
\begin{frame}{Model}{The Household's Dynamic Problem}\label{fr:hh}
  \small
  The representative household in country $i$ chooses consumption, investment,
  capital, and bonds to maximize:
  \begin{equation*}
    \max \sum_{t=0}^{\infty} \delta^t L_{i,t} \ln \left(\bar{C}_{i,t}\right)
  \end{equation*}

  \vspace{4pt}
  subject to
  \begin{align*}
    E^C_{i,t} + \frac{E^X_{i,t}}{\dev{\xi_{i,t}}} + B_{i,t+1}
      + \dev{\frac{\phi_{i,t}}{2}\left(B_{i,t+1}-\bar{B}_i\right)^2}
      &= w_{i,t} L_{i,t} + \frac{r_{i,t} K_{i,t}}{\dev{\xi_{i,t}}} + R_t B_{i,t} + T_{i,t} \\
    K_{i,t+1} &= (1-d) K_{i,t} + \left(X_{i,t}\right)^\sigma \left(K_{i,t}\right)^{1-\sigma}
  \end{align*}

  \vspace{4pt}
  where $\bar{C}_{i,t}$ is defined implicitly (CLM, 2017):
  \begin{equation*}
    \sum_{j \in \mathcal{J}} \left(\mu^j_{i,t}\right)^{\frac{1}{\psi}}
    \left(\frac{C_{i,t}}{L_{i,t}}\right)^{\dev{\frac{1-\psi}{\psi}\varepsilon^j}}
    \left(\frac{C^j_{i,t}}{L_{i,t}}\right)^{\frac{\psi-1}{\psi}} = 1
  \end{equation*}
  and $\psi > 0$, $\mu^j_{i,t} > 0$, and $\sum_j \varepsilon^j = 1$.\ap{ap:cal}

  {\footnotesize \dev{Coloured}: the three departures from Reyes-Heroles (2018) --- a
  capital wedge, a portfolio adjustment cost in place of a preference shifter, and the CLM
  parameterization of the aggregator.}
\end{frame}
\begin{frame}{Model}{Methodology}\label{fr:meth}
  \vfill
  \linespread{1.3}\selectfont
  \begin{enumerate}\itemsep12pt
    \item Invert the model and find shocks.
    \item Given parameters, $K_{i,1}$, $B_{i,1}=0$, and paths for
          $\{\mu^j_{i,t},T^j_{i,t},\tau^j_{ih,t},\xi_{i,t},\phi_{i,t}\}$, find
          paths of prices, quantities, trade shares and bonds such that all
          equilibrium conditions hold.\ap{ap:ek}
    \item Compare results against counterfactual equilibria with fixed trade
          costs and financial frictions.
  \end{enumerate}
  \vfill
\end{frame}
\begin{frame}[fragile]{Model}{Equilibrium Conditions\ap{ap:mc}}
  \centering
  \tiny
  % The table runs a shade past the frame at \tiny, so scale it rather than clip
  % the last row.
  \IfFileExists{tables/equilibrium_all.tex}
    {\scalebox{0.92}{% Save this file as: tables/equilibrium_all.tex
\begin{tabular}{p{0.5cm} p{7.0cm} p{2.0cm}}
\toprule
\textbf{Tag} & \textbf{Equation / Recovery Formula} & \textbf{Description} \\
\midrule
(C1) & $\displaystyle P^C_{i,t} \equiv \left[\sum_j \mu^j_{i,t}\left(\frac{C_{i,t}}{L_{i,t}}\right)^{(1-\psi)(\varepsilon^j-1)}(P^j_{i,t})^{1-\psi}\right]^{\frac{1}{1-\psi}}$ & Consumption price index \\

(X1) & $\displaystyle P^X_{i,t} \equiv \left[\sum_{j} \upsilon^j_{i}(P^j_{i,t})^{1-\zeta}\right]^{\frac{1}{1-\zeta}}$ & Investment price index \\

(F1) & $\displaystyle c^j_{i,t} = \kappa^j_i \left[(r_{i,t}^{\varphi_i} w_{i,t}^{1-\varphi_i})^{\beta^j_i} \prod_m (P^m_{i,t})^{\nu^{j,m}_i}\right]$ & Sectoral unit cost \\

(T1) & $\displaystyle \pi^j_{ih,t} = \frac{T^j_{h,t}(c^j_{h,t}\tau^j_{ih,t})^{-\theta}}{\Phi^j_{i,t}}$ & Trade share (Gravity) \\

(E1) & $\displaystyle \frac{E^C_{i,t+1}}{E^C_{i,t}} = \frac{L_{i,t+1}}{L_{i,t}}\frac{\widetilde{\phi}_{i,t+1}}{\widetilde{\phi}_{i,t}}\frac{\delta R_{t+1}}{1+\phi_{i,t}(B_{i,t+1}-\bar{B}_i)}$ & Bond Euler equation \\

(K1) & $\displaystyle \frac{R_{t+1}}{1+\phi_{i,t}(B_{i,t+1}-\bar{B}_i)} P^X_{i,t} = \frac{\xi_{i,t}}{\xi_{i,t+1}}\left[r_{i,t+1} + (1-d)P^X_{i,t+1}\right]$ & Capital Euler ($\sigma=1$) \\
\midrule
(INV1) & $\displaystyle T^j_{i,t} = \pi^j_{ii,t}\left(\Gamma\,\frac{c^j_{i,t}}{P^j_{i,t}}\right)^{\theta}$ & Productivity inversion \\

(INV2) & $\displaystyle \tau^j_{ih,t} = \frac{P^j_{i,t}}{P^j_{h,t}} \left(\frac{\pi^j_{hh,t}}{\pi^j_{ih,t}}\right)^{\frac{1}{\theta}}$ & Trade cost inversion \\

(INV3) & $\displaystyle \phi_{R,t} = \frac{1}{B_{R,t+1}-\bar{B}_{R}} \left[ \frac{E^C_{R,t}}{E^C_{R,t+1}}\frac{L_{R,t+1}}{L_{R,t}}\frac{\widetilde{\phi}_{R,t+1}}{\widetilde{\phi}_{R,t}}\delta R_{t+1} - 1 \right]$ & Financial friction inversion \\

(INV4) & $\displaystyle \frac{\xi_{i,t}}{\xi_{i,t+1}} = R_{t+1}P^X_{i,t}\Big/\left[1+\phi_{i,t}(B_{i,t+1}-\bar{B}_i)\right]\left[r_{i,t+1}+(1-d)P^X_{i,t+1}\right]$ & Capital wedge inversion \\
\bottomrule
\end{tabular}
}}
    {\normalsize\textit{[tables/equilibrium\_all.tex not found]}}
\end{frame}
%=======================================================
\section{Recovering parameters}
%=======================================================
\begin{frame}{Recovering parameters}{Data}\label{fr:data}
  \small
  One two-country, three-sector panel, 1965--2014: the U.S.\ against an aggregated
  rest of the world,\ap{ap:row} in agriculture, manufacturing and services.
  \vspace{6pt}

  \centering \footnotesize
  \setlength{\tabcolsep}{5pt}
  \begin{tabular}{@{}L{3.2cm}lL{4.4cm}@{}}
    \toprule
    Source & Years & Contribution \\
    \midrule
    Long-Run WIOD, v1.1
      & 1965--2000
      & Input--output tables, gross output, value added, deflators \\[2pt]
    WIOD 2016 release
      & 2000--2014
      & The same, at finer country and sector detail \\[2pt]
    GGDC Productivity Level DB
      & 2005
      & Sectoral PPP benchmark\ap{ap:prices} \\[2pt]
    Penn World Table 11.0
      & 1965--2014
      & Employment, capital, investment prices \\
    \bottomrule
  \end{tabular}

  \vspace{8pt}
  \raggedright \small
  \structure{\textbf{Two constructions do the work.}} The releases are chain-linked at
  2000, levels through 2000 and growth rates after. Sectoral prices are then put in
  common units by anchoring each country's deflator to the 2005 PPP benchmark --- which
  is what lets trade costs and productivities be separated.
\end{frame}
\figframe{Recovering parameters}{Consumption Demand Shifters}{consumption_demand_shifters.pdf}
\figframe{Recovering parameters}{Trade Costs\ap{ap:row}}{trade_costs.pdf}
\figframe{Recovering parameters}{Productivities}{productivities.pdf}
\figframe{Recovering parameters}{Financial Friction}{financial_friction.pdf}
\figframe{Recovering parameters}{Investment Efficiency}{investment_efficiency.pdf}
\begin{frame}[plain]
  \vfill
  \centering
  {\usebeamercolor[fg]{structure}\Large Thank you}
  \vfill
\end{frame}
%=======================================================
\section*{Appendix}
%=======================================================
\begin{frame}{Appendix}{Sectoral Prices and Trade Shares (Eaton--Kortum)}\label{ap:ek}
  \begin{itemize}\itms
    \item \structure{\textbf{Ideal price index.}} Sectoral prices in each $j$ are given by
      \begin{equation*}
        P^j_{i,t} = \Gamma \left(\Phi^j_{i,t}\right)^{-\frac{1}{\theta}},
        \quad \text{where} \quad
        \Phi^j_{i,t} = \sum_{h \in \mathcal{I}} T^j_{h,t}
        \left(c^j_{h,t} \tau^j_{ih,t}\right)^{-\theta}
      \end{equation*}
    \item \structure{\textbf{Unit cost.}} The minimum cost of the input bundle $c^j_{i,t}$ is
      \begin{equation*}
        c^j_{i,t} = \kappa^j_i
        \left[ \left(r_{i,t}\right)^{\varphi_i} \left(w_{i,t}\right)^{1-\varphi_i} \right]^{\beta^j_i}
        \left[ \prod_{k \in \mathcal{J}} \left(P^k_{i,t}\right)^{\nu^{j,k}_i} \right]^{1-\beta^j_i}
      \end{equation*}
    \item \structure{\textbf{Trade shares.}} Define $E^j_{i,t} \equiv P^j_{i,t} Q^j_{i,t}$ as
      total expenditure on sector $j$. The share of expenditure on goods imported from $h$ is
      \begin{equation*}
        \pi^j_{ih,t} \equiv \frac{E^j_{ih,t}}{E^j_{i,t}}
        = \frac{T^j_{h,t} \left(c^j_{h,t} \tau^j_{ih,t}\right)^{-\theta}}{\Phi^j_{i,t}}
      \end{equation*}
      the multisector version of the gravity equation.
  \end{itemize}
  \apback{fr:meth}
\end{frame}
\begin{frame}{Appendix}{Sectoral Price Levels: Construction}\label{ap:prices}
  \footnotesize
  The goal is a price for each country, sector and year that is comparable
  \emph{across} countries, not only over time.
  \begin{enumerate}\itemsep2pt
    \item \structure{\textbf{Deflators.}} Gross output price indices from both WIOD releases.
          The long-run release is in USD, so it is multiplied by the
          \emph{2005} exchange rate to put it in local currency; the 2016
          release is already in local currency.
    \item \structure{\textbf{Splice.}} Chain-link at 2000 --- long-run levels through 2000,
          2016 growth rates after --- then normalize to 2005 within each
          country and sector, so $\tilde{\pi}_{i,c,2005}=1$.
    \item \structure{\textbf{Benchmark.}} Take the 2005 gross-output PPP for each country
          and sector from PLD, $\text{PPP}_{i,c}$, in local currency per USD.
          It is a gross-output-weighted geometric mean of the product-level PPPs.
    \item \structure{\textbf{Combine,}} holding the exchange rate at its 2005 value:
      \begin{equation*}
        P_{i,c,t} = \tilde{\pi}_{i,c,t}\cdot\frac{\text{PPP}_{i,c}}{e_{c,2005}}
      \end{equation*}
  \end{enumerate}
  At 2005 the level is the sector's PPP relative to the United States; in every
  other year it moves only with that country's own deflator. Fixing the exchange
  rate keeps the nominal dollar cycle out of the cross-country comparison. For
  the United States $\text{PPP}=e=1$, so the series is a pure within-country
  price index.
  \apback{fr:data}
\end{frame}
\begin{frame}{Appendix}{Sectoral Price Levels: Aggregating to ROW}\label{ap:row}
  \small
  The model has two countries, so member price levels have to be combined into
  one. Gross-output-weighted geometric mean of the \emph{levels}, computed
  within each year:
  \begin{equation*}
    \ln P_{i,\text{ROW},t}=\sum_{c\in\text{ROW}}w_{c,t}\ln P_{i,c,t},
    \qquad
    w_{c,t}=\frac{Y_{c,t}}{\sum_{c'\in\text{ROW}}Y_{c',t}}
  \end{equation*}
  \begin{itemize}\itms
    \item A mean of levels, not a chained index of changes: a cheap member
          lowers the ROW price level as soon as it gains share.
    \item A chained index reads only price \emph{changes}, so an entrant
          contributes nothing on arrival and its low level is never picked up.
          Using one drove measured ROW services productivity negative.
    \item Everything is aggregated to two countries first, and the disturbances
          are recovered from the aggregated data, so no member-level parameter
          is pooled. Bilateral trade costs follow directly from the aggregated
          prices and trade shares,
          $\tau^j_{ih,t}=\dfrac{P^j_{i,t}}{P^j_{h,t}}
           \left(\dfrac{\pi^j_{hh,t}}{\pi^j_{ih,t}}\right)^{1/\theta}$.
  \end{itemize}
  \apback{fr:data}
\end{frame}
\begin{frame}{Appendix}{Market Clearing Conditions}\label{ap:mc}
  \begin{itemize}\itms
    \item \structure{\textbf{Goods and factor markets.}}
      \begin{equation*}
        C^j_{i,t} + X^j_{i,t} + \sum_{k \in \mathcal{J}} D^{k,j}_{i,t} = Q^j_{i,t}
        \quad \text{for all } j \in \mathcal{J}
      \end{equation*}
      \begin{equation*}
        \sum_{j \in \mathcal{J}} L^j_{i,t} = L_{i,t} \quad \text{and} \quad
        \sum_{j \in \mathcal{J}} K^j_{i,t} = K_{i,t}
      \end{equation*}
    \item \structure{\textbf{Value of production.}} Let $Y^j_{i,t}$ denote the gross value of
      production.
      \begin{equation*}
        Y^j_{i,t} = \sum_{h \in \mathcal{I}} \pi^j_{hi,t} E^j_{h,t}
        \quad \text{for all } j \in \mathcal{J}
      \end{equation*}
    \item \structure{\textbf{Country resource constraint.}}
      \begin{equation*}
        B_{i,t+1} - R_t B_{i,t} = \sum_{j \in \mathcal{J}} Y^j_{i,t} - E^j_{i,t}
      \end{equation*}
    \item \structure{\textbf{International financial markets clear.}}
      \begin{equation*}
        \sum_{i \in \mathcal{I}} B_{i,t+1} = 0 \quad \text{for all } t
      \end{equation*}
  \end{itemize}
  \apback{fr:meth}
\end{frame}
\begin{frame}{Appendix}{Calibration}\label{ap:cal}
  \begin{center}
  \begin{tabular}{@{}llc@{}}
    \toprule
    $\theta$   & trade elasticity              & 4 \\
    $\eta$     & elasticity across varieties   & 2 \\
    $\psi$     & elasticity across sectors     & 0.57 \\
    $\zeta$    & investment elasticity         & 1 \\
    $\sigma$   & capital adjustment            & 0 \\
    $\delta$   & discount factor               & 0.95 \\
    $d$        & depreciation                  & 0.05 \\
    \bottomrule
  \end{tabular}
  \end{center}
  \vspace{8pt}
  $\varepsilon^j$ from Comin, Lashkari and Mestieri (2021); $\varphi_i$ from
  Caselli and Feyrer (2007); $\beta^j_i$ and $\nu^{j,m}_i$ are sample averages.
  \apback{fr:hh}
\end{frame}
\end{document}
